% vim: set ft=tex tabstop=4 shiftwidth=4 noexpandtab:

\input{../../include/latex/tikz/header.tex}

% opening %{{{1

\begin{document}
\begin{tikzpicture}[scale=1.0]

% parameters %{{{1

	\tikzmath{
		\radius = 5.4 ;
		\rotation = 90 ;
		\numsidesfloat = 5.0 ;
		\numsides = int(\numsidesfloat) ;
		int \i ;
		for \i in {1,...,\numsides} {
			\ang{\i} = \rotation + 360 / \numsidesfloat * (\i-1) ;
		};
		\xmin = - \radius - 1 ;
		\xmax = \radius + 1 ;
		\ymin = - \radius / 6 ;
		\ymax = \radius + 1 ;
	}

% coordinates %{{{1

	\coordinate (O) at (0,0);

	\foreach \i in {1,...,\numsides}
		\coordinate (P_\i) at (\ang{\i}:\radius);

% intersections %{{{1

	\begin{pgfinterruptboundingbox}
		\foreach \i
			[
				evaluate=\i as \j using { int(mod(\i, \numsides) + 1) },
				evaluate=\i as \k using { int(mod(\i + 1, \numsides) + 1) },
			]
			in {1,...,\numsides}
		{
			\tkzInterLL(O,P_\j)(P_\i,P_\k)\tkzGetPoint{I_\j}
		}
	\end{pgfinterruptboundingbox}

	\begin{pgfinterruptboundingbox}
		\foreach \i
			[
				evaluate=\i as \j using { int(mod(\i, \numsides) + 1) },
				evaluate=\i as \k using { int(mod(\i + 1, \numsides) + 1) },
				evaluate=\i as \l using { int(mod(\i + 2, \numsides) + 1) }
			]
			in {1,...,\numsides}
		{
			\tkzInterLL(P_\i,P_\k)(P_\j,P_\l)\tkzGetPoint{S_\j}
		}
	\end{pgfinterruptboundingbox}

% points, dots, vertices %{{{1

	\fill (O) circle (0.4mm);

	\foreach \i in {1,...,\numsides}
		\fill (P_\i) circle (0.4mm);

	\foreach \i in {1,...,\numsides}
		\fill (I_\i) circle (0.4mm);

	\foreach \i in {1,...,\numsides}
		\fill (S_\i) circle (0.4mm);

% segments, sides, lines %{{{1

	\draw (P_1) \foreach \i in {2,...,\numsides} { -- (P_\i) } -- cycle;

	\draw (P_1) -- (P_3) -- (P_5) -- (P_2) -- (P_4) -- cycle;

	\foreach \i in {1,...,\numsides}
		\draw[dashed] (O) -- (P_\i);

	\foreach \i in {1,...,\numsides}
		\draw[dashed] (O) -- (S_\i);

% circles %{{{1

	\draw (O) circle (\radius);

% points, dots, vertices labels %{{{1

	\node[anchor=east] at ($ (O) + (250:7pt) $) {$O$};

	\node[above] at (P_1) {$P_1$};
	\node[left] at (P_2) {$P_2$};
	\node[below left] at (P_3) {$P_3$};
	\node[below right] at (P_4) {$P_4$};
	\node[right] at (P_5) {$P_5$};

	\node[above left] at (I_1) {$I_1$};
	\node[left] at ($ (I_2) + (240:3pt) $) {$I_2$};
	\node[below] at (I_3) {$I_3$};
	\node[right] at (I_4) {$I_4$};
	\node[above right] at (I_5) {$I_5$};

	\node[above left] at (S_1) {$S_1$};
	\node[left] at (S_2) {$S_2$};
	\node[below] at (S_3) {$S_3$};
	\node[right] at (S_4) {$S_4$};
	\node[above right] at (S_5) {$S_5$};

% segments, sides, lines labels %{{{1

	%\node[above left] at ($ (P_1)!0.65!(S_1) $) {$a$};

% segments, sides, lines marks %{{{1

	\foreach \i in {1,...,\numsides}
		\tkzMarkSegments[mark=||, size=2pt, pos=0.65](P_\i,S_\i);

	\foreach \i
		[
			evaluate=\i as \j using { int(mod(\i,\numsides) + 1) }
		]
		in {1,...,\numsides}
	{
		\tkzMarkSegments[mark=||, size=2pt, pos=0.65](P_\j,S_\i);
	}

	\foreach \i in {1,...,\numsides}
		\tkzMarkSegments[mark=|, size=2pt, pos=0.5](I_\i,S_\i);

	\foreach \i
		[
			evaluate=\i as \j using { int(mod(\i,\numsides) + 1) }
		]
		in {1,...,\numsides}
	{
		\tkzMarkSegments[mark=|, size=2pt, pos=0.5](I_\j,S_\i);
	}

	\foreach \i in {1,...,\numsides} {
		\tkzMarkSegments[mark=|, size=2pt, pos=0.65](P_\i,I_\i)
		\tkzMarkSegments[mark=|, size=2pt, pos=0.70](P_\i,I_\i)
		\tkzMarkSegments[mark=|, size=2pt, pos=0.75](P_\i,I_\i)
	}

	\foreach \i in {1,...,\numsides} {
		\tkzMarkSegments[mark=||, size=2pt, pos=0.65](O,S_\i)
		\tkzMarkSegments[mark=||, size=2pt, pos=0.75](O,S_\i)
	}

	\foreach \i in {1,...,\numsides} {
		\tkzMarkSegments[mark=||, size=2pt, pos=0.65](O,I_\i)
		\tkzMarkSegments[mark=||, size=2pt, pos=0.70](O,I_\i)
		\tkzMarkSegments[mark=||, size=2pt, pos=0.75](O,I_\i)
	}

% circles labels %{{{1

	\node at ($ (O) + (60:\radius + 0.3) $) {$\mathscr{C}$};

% angles labels %{{{1

	% ------- between radius and diagonal

	\foreach \i
		[
			evaluate=\i as \j using { int(mod(\i - 1 + 2, \numsides) + 1) }
		]
		in {1,...,\numsides}
	{
		\pic[draw, ->, "$18^\circ$", angle radius=1.5cm, angle eccentricity=1.2]
			{angle = P_\j--P_\i--O};
	}

	\foreach \i
		[
			evaluate=\i as \j using { int(mod(\i - 1 + 3, \numsides) + 1) }
		]
		in {1,...,\numsides}
	{
		\pic[draw, ->, "$18^\circ$", angle radius=1.5cm, angle eccentricity=1.2]
			{angle = O--P_\i--P_\j};
	}

	% ------- isosceles triangles outside the pentagram

	\foreach \i
		[
			evaluate=\i as \j using { int(mod(\i - 1 + 1, \numsides) + 1) },
			evaluate=\i as \k using { int(mod(\i - 1 + 2, \numsides) + 1) }
		]
		in {1,...,\numsides}
	{
		\pic[draw, ->, "$36^\circ$", angle radius=1.1cm, angle eccentricity=1.3]
			{angle = P_\i--P_\k--P_\j};
		\pic[draw, ->, "$36^\circ$", angle radius=1.1cm, angle eccentricity=1.3]
			{angle = P_\j--P_\i--P_\k};
	}

	% ------- around center

	\foreach \i in {1,...,\numsides}
		\pic[draw, ->, "$36^\circ$", angle radius=0.8cm, angle eccentricity=1.4]
			{angle = I_\i--O--S_\i};

	\foreach \i
		[
			evaluate=\i as \j using { int(mod(\i - 1 + 1, \numsides) + 1) }
		]
		in {1,...,\numsides}
	{
		\pic[draw, ->, "$36^\circ$", angle radius=0.8cm, angle eccentricity=1.4]
			{angle = S_\i--O--I_\j};
	}

% right angles markers %{{{2

	\foreach \i
		[
			evaluate=\i as \j using { int(mod(\i, \numsides) + 1) }
		]
		in {1,...,\numsides}
	{
		\pic [draw, angle radius=7pt, angle eccentricity=1]
			{right angle = P_\i--I_\j--P_\j};
	}

% closing %{{{1

\end{tikzpicture}
\end{document}
